% ========================================================
% BÀI TẬP TOÁN 5 - KẾT NỐI TRI THỨC VỚI CUỘC SỐNG
% Bài 16: Các đơn vị đo diện tích
% Định dạng: BTTeX Part III - Trắc nghiệm điền số ngắn
% ========================================================

% Câu 1
\begin{ex}
  Một mét vuông bằng bao nhiêu đề-xi-mét vuông?
  \shortans{100}
  \loigiai{
    $1\text{ m}^2 = 100\text{ dm}^2$.
  }
\end{ex}

% Câu 2
\begin{ex}
  Một đề-xi-mét vuông bằng bao nhiêu xăng-ti-mét vuông?
  \shortans{100}
  \loigiai{
    $1\text{ dm}^2 = 100\text{ cm}^2$.
  }
\end{ex}

% Câu 3
\begin{ex}
  Một xăng-ti-mét vuông bằng bao nhiêu mi-li-mét vuông?
  \shortans{100}
  \loigiai{
    $1\text{ cm}^2 = 100\text{ mm}^2$.
  }
\end{ex}

% Câu 4
\begin{ex}
  Một mét vuông bằng bao nhiêu xăng-ti-mét vuông?
  \shortans{10000}
  \loigiai{
    $1\text{ m}^2 = 100 \times 100 = 10\,000\text{ cm}^2$.
  }
\end{ex}

% Câu 5
\begin{ex}
  Đổi $6\text{ m}^2\text{ } 45\text{ dm}^2$ thành bao nhiêu đề-xi-mét vuông?
  \shortans{645}
  \loigiai{
    $6\text{ m}^2 = 600\text{ dm}^2$. Do đó $6\text{ m}^2\text{ } 45\text{ dm}^2 = 645\text{ dm}^2$.
  }
\end{ex}

% Câu 6
\begin{ex}
  Đổi $8\text{ dm}^2\text{ } 5\text{ cm}^2$ thành bao nhiêu xăng-ti-mét vuông?
  \shortans{805}
  \loigiai{
    $8\text{ dm}^2 = 800\text{ cm}^2$. Do đó $8\text{ dm}^2\text{ } 5\text{ cm}^2 = 805\text{ cm}^2$.
  }
\end{ex}

% Câu 7
\begin{ex}
  Đổi $3\text{ cm}^2\text{ } 20\text{ mm}^2$ thành bao nhiêu mi-li-mét vuông?
  \shortans{320}
  \loigiai{
    $3\text{ cm}^2 = 300\text{ mm}^2$. Do đó $3\text{ cm}^2\text{ } 20\text{ mm}^2 = 320\text{ mm}^2$.
  }
\end{ex}

% Câu 8
\begin{ex}
  Đổi $500\text{ dm}^2$ thành bao nhiêu mét vuông?
  \shortans{5}
  \loigiai{
    $500 : 100 = 5\text{ m}^2$.
  }
\end{ex}

% Câu 9
\begin{ex}
  Đổi $12\,000\text{ cm}^2$ thành bao nhiêu đề-xi-mét vuông?
  \shortans{120}
  \loigiai{
    $12\,000 : 100 = 120\text{ dm}^2$.
  }
\end{ex}

% Câu 10
\begin{ex}
  Một viên gạch men hình vuông có cạnh dài $50\text{ cm}$. Diện tích một viên gạch là bao nhiêu xăng-ti-mét vuông?
  \shortans{2500}
  \loigiai{
    Diện tích viên gạch là $50 \times 50 = 2\,500\text{ cm}^2$.
  }
\end{ex}

% Câu 11
\begin{ex}
  Một phòng học có sàn hình chữ nhật diện tích $60\text{ m}^2$. Người ta lát nền bằng gạch có diện tích mỗi viên là $25\text{ dm}^2$. Cần bao nhiêu viên gạch (bỏ qua mạch vữa)?
  \shortans{240}
  \loigiai{
    $60\text{ m}^2 = 6\,000\text{ dm}^2$. Số viên gạch cần là: $6\,000 : 25 = 240$ (viên).
  }
\end{ex}

% Câu 12
\begin{ex}
  Một tấm kính hình chữ nhật dài $12\text{ dm}$, rộng $8\text{ dm}$. Diện tích tấm kính là bao nhiêu đề-xi-mét vuông?
  \shortans{96}
  \loigiai{
    Diện tích là: $12 \times 8 = 96\text{ dm}^2$.
  }
\end{ex}

% Câu 13
\begin{ex}
  Một con tem thư hình chữ nhật có kích thước $3\text{ cm} \times 2\text{ cm}$. Diện tích con tem thư là bao nhiêu mi-li-mét vuông?
  \shortans{600}
  \loigiai{
    Diện tích là $3 \times 2 = 6\text{ cm}^2 = 600\text{ mm}^2$.
  }
\end{ex}

% Câu 14
\begin{ex}
  Đổi $4\text{ m}^2\text{ } 800\text{ cm}^2$ thành bao nhiêu đề-xi-mét vuông?
  \shortans{408}
  \loigiai{
    $4\text{ m}^2 = 400\text{ dm}^2$; $800\text{ cm}^2 = 8\text{ dm}^2$. Tổng là $400 + 8 = 408\text{ dm}^2$.
  }
\end{ex}

% Câu 15
\begin{ex}
  Người ta làm một khung ảnh hình vuông cạnh $30\text{ cm}$. Diện tích khung ảnh đó là bao nhiêu xăng-ti-mét vuông?
  \shortans{900}
  \loigiai{
    Diện tích là $30 \times 30 = 900\text{ cm}^2$.
  }
\end{ex}

% Câu 16
\begin{ex}
  Một bức tường hình chữ nhật dài $6\text{ m}$, cao $3\text{ m}$. Diện tích bức tường là bao nhiêu mét vuông?
  \shortans{18}
  \loigiai{
    Diện tích bức tường là $6 \times 3 = 18\text{ m}^2$.
  }
\end{ex}

% Câu 17
\begin{ex}
  Để sơn $1\text{ m}^2$ tường cần hết $200\text{ g}$ sơn. Để sơn bức tường diện tích $25\text{ m}^2$ cần bao nhiêu ki-lô-gam sơn?
  \shortans{5}
  \loigiai{
    Lượng sơn cần là: $25 \times 200 = 5\,000\text{ g} = 5\text{ kg}$.
  }
\end{ex}

% Câu 18
\begin{ex}
  Đổi $2\text{ dam}^2$ thành bao nhiêu mét vuông?
  \shortans{200}
  \loigiai{
    $1\text{ dam}^2 = 100\text{ m}^2$. Do đó $2\text{ dam}^2 = 200\text{ m}^2$.
  }
\end{ex}

% Câu 19
\begin{ex}
  Đổi $15\text{ hm}^2$ thành bao nhiêu héc-ta?
  \shortans{15}
  \loigiai{
    $1\text{ hm}^2 = 1\text{ ha}$. Do đó $15\text{ hm}^2 = 15\text{ ha}$.
  }
\end{ex}

% Câu 20
\begin{ex}
  Đổi $3\text{ dam}^2\text{ } 20\text{ m}^2$ thành bao nhiêu mét vuông?
  \shortans{320}
  \loigiai{
    $3\text{ dam}^2 = 300\text{ m}^2$. Do đó $3\text{ dam}^2\text{ } 20\text{ m}^2 = 320\text{ m}^2$.
  }
\end{ex}

% Câu 21
\begin{ex}
  Một mảnh vườn hình bình hành có đáy $15\text{ m}$ và chiều cao $8\text{ m}$. Diện tích mảnh vườn là bao nhiêu mét vuông?
  \shortans{120}
  \loigiai{
    Diện tích hình bình hành là: $15 \times 8 = 120\text{ m}^2$.
  }
\end{ex}

% Câu 22
\begin{ex}
  Một bảng mica hình thoi có hai đường chéo dài $80\text{ cm}$ và $50\text{ cm}$. Diện tích bảng mica là bao nhiêu đề-xi-mét vuông?
  \shortans{20}
  \loigiai{
    Diện tích là $(80 \times 50) : 2 = 2\,000\text{ cm}^2 = 20\text{ dm}^2$.
  }
\end{ex}

% Câu 23
\begin{ex}
  Một tấm bảng thông báo hình chữ nhật dài $2\text{ m}$, rộng $1\text{ m } 20\text{ cm}$. Diện tích tấm bảng là bao nhiêu đề-xi-mét vuông?
  \shortans{240}
  \loigiai{
    $2\text{ m} = 20\text{ dm}$; $1\text{ m } 20\text{ cm} = 12\text{ dm}$. Diện tích là $20 \times 12 = 240\text{ dm}^2$.
  }
\end{ex}

% Câu 24
\begin{ex}
  Nếu gấp độ dài cạnh của một hình vuông lên $3$ lần thì diện tích hình vuông gấp lên mấy lần?
  \shortans{9}
  \loigiai{
    Diện tích hình vuông tỉ lệ với bình phương cạnh: $3 \times 3 = 9$ (lần).
  }
\end{ex}

% Câu 25
\begin{ex}
  Nếu giảm chiều dài hình chữ nhật đi một nửa và gấp đôi chiều rộng thì diện tích hình chữ nhật thay đổi gấp mấy lần?
  \shortans{1}
  \loigiai{
    Diện tích mới bằng $\frac{1}{2} \times 2 = 1$ lần diện tích cũ (không thay đổi).
  }
\end{ex}
